\clearpage


\chapter{Технологическая часть}\label{ch:-2}

В этом разделе обоснованы средства реализации, а так же представлены реализации алгоритмов и функциональные тесты.


\section{Средства реализации}\label{sec:-2}

Для реализации алгоритмов в этой лабораторной работы был выбран язык $Python^{[3]}$ по следующим причинам:
\begin{itemize}
    \item python предоставляет необходимые структуры данных, такие как списки (массивы) и вложенные списки (матрицы);
    \item поддерживаемые библиотеки, такие как $matplotlib{[4]}$ позволяют визуализировать результаты исследования.
\end{itemize}

\clearpage


\section{Реализация алгоритмов}

В листингах~\ref{lst:brf} ---~\ref{lst:ant} представлены реализации алгоритмов.

\begin{lstlisting}[style=python, label=lst:brf,caption=Алгоритм полного перебора]
def generate_permutations(arr):
    result = []
    def permute(a, l):
        if l == len(a):
            result.append(a[:])
        else:
            for i in range(l, len(a)):
                a[l], a[i] = a[i], a[l]
                permute(a, l+1)
                a[l], a[i] = a[i], a[l]
    permute(arr, 0)
    return result


def tsp_brute_force(distance_matrix):
    n = len(distance_matrix)
    vertices = list(range(n))

    all_permutations = generate_permutations(vertices)

    best_length = float('inf')
    best_route = None

    for route in all_permutations:
        route_length = 0.0
        for i in range(len(route)-1):
            route_length += distance_matrix[route[i]][route[i+1]]
        if route_length < best_length:
            best_length = route_length
            best_route = route[:]

    return best_route, best_length
\end{lstlisting}
\clearpage


\begin{lstlisting}[style=python, label=lst:ant,caption=Муравьиный алгоритм]
def ant_colony_tsp_with_resources(distance_matrix,
                                  num_ants=10,
                                  num_iterations=100,
                                  alpha=0.5,
                                  beta=0.5,
                                  rho=0.5,
                                  tau_0=0.01,
                                  max_resources=100.0):
    n = len(distance_matrix)

    total_dist = 0.0
    for i in range(n):
        for j in range(n):
            if i != j:
                total_dist += distance_matrix[i][j]
    Q = total_dist / n

    pheromone = [[tau_0 for _ in range(n)] for _ in range(n)]
    eta = [[0 if distance_matrix[i][j] == 0 else 1.0 / distance_matrix[i][j] for j in range(n)] for i in range(n)]

    best_route = None
    best_length = float('inf')

    for iteration in range(num_iterations):
        all_routes = []
        all_lengths = []

        for ant in range(num_ants):
            start = ant % n
            visited = [False] * n
            visited[start] = True
            route = [start]
            current = start

            remaining_resources = max_resources

            while len(route) < n:

                candidates = []
                for j in range(n):
                    if not visited[j]:
                        cost = distance_matrix[current][j]
                        if remaining_resources -- cost >= 0:
                            val = (pheromone[current][j] ** alpha) * (eta[current][j] ** beta)
                            candidates.append((j, val))

                if len(candidates) == 0:
                    break

                denom = sum(val for _, val in candidates)
                r = random.random()
                cumulative = 0.0
                next_city = None
                for (city, val) in candidates:
                    p = val / denom
                    cumulative += p
                    if r <= cumulative:
                        next_city = city
                        break

                visited[next_city] = True
                route.append(next_city)

                cost = distance_matrix[current][next_city]
                remaining_resources -= cost
                current = next_city

            route_length = 0.0
            for i in range(len(route) - 1):
                route_length += distance_matrix[route[i]][route[i + 1]]

            all_routes.append(route)
            all_lengths.append(route_length)

            if route_length < best_length:
                best_length = route_length
                best_route = route[:]


        for i in range(n):
            for j in range(n):
                pheromone[i][j] *= (1 - rho)
                if pheromone[i][j] < 1e-12:
                    pheromone[i][j] = 1e-12

        for k, route in enumerate(all_routes):
            route_length = all_lengths[k]
            if route_length > 0:
                delta = Q / route_length
            else:
                delta = Q
            for i in range(len(route) - 1):
                a, b = route[i], route[i + 1]
                pheromone[a][b] += delta
                pheromone[b][a] += delta

    return best_route, best_length
\end{lstlisting}


\section{Тестирование}

При тестирование муравьиного алгоритма нельзя проверить при помощи сравнения с точным значением, в отличие от полного
перебора. Можно лишь убедиться что результирующий путь содержит все вершины.

\subsubsection*{Тест 1}

\textbf{Входные данные:}
\[
    \begin{bmatrix}
        0  & 2 & 9  & 10 & 5  \\
        1  & 0 & 6  & 4  & 7  \\
        15 & 7 & 0  & 8  & 3  \\
        6  & 3 & 12 & 0  & 11 \\
        9  & 5 & 2  & 6  & 0  \\
    \end{bmatrix}
\]

\textbf{Ожидаемый лучший маршрут:} [3, 1, 0, 4, 2], длина = 11


\textbf{Выходные данные полного перебора:} [3, 1, 0, 4, 2], длина = 11


\textbf{Выходные данные муравьиного алгоритма:} [3, 1, 0, 4, 2], длина = 11

\subsubsection{Тест 2}

\textbf{Входные данные:}
\[
    \begin{bmatrix}
        0  & 8  & 4  & 1  & 9  & 27 & 5  & 8  & 12 & 14 \\
        8  & 0  & 13 & 4  & 14 & 12 & 13 & 7  & 20 & 4  \\
        4  & 13 & 0  & 5  & 16 & 8  & 18 & 4  & 21 & 17 \\
        1  & 4  & 5  & 0  & 7  & 12 & 9  & 8  & 13 & 16 \\
        9  & 14 & 16 & 7  & 0  & 12 & 17 & 12 & 17 & 4  \\
        27 & 12 & 8  & 12 & 12 & 0  & 8  & 13 & 9  & 4  \\
        5  & 13 & 18 & 9  & 17 & 8  & 0  & 10 & 13 & 1  \\
        8  & 7  & 4  & 8  & 12 & 13 & 10 & 0  & 7  & 11 \\
        12 & 20 & 21 & 13 & 17 & 9  & 13 & 7  & 0  & 10 \\
        14 & 4  & 17 & 16 & 4  & 4  & 1  & 11 & 10 & 0  \\
    \end{bmatrix}
\]


\textbf{Ожидаемый лучший маршрут:} [1, 3, 0, 2, 7, 8, 5, 6, 9, 4], длина = 42


\textbf{Выходные данные полного перебора:} [1, 3, 0, 2, 7, 8, 5, 6, 9, 4], длина = 42


\textbf{Выходные данные муравьиного алгоритма:} [4, 9, 6, 0, 3, 1, 7, 2, 5, 8], длина = 43

\section*{Вывод}

В технологической части были представлены листинги реализованных алгоритмов,
включая алгоритм полного перебора и муравьиный алгоритм.

